\begin{textsample}{NEG Dim 6 – global\_south\_2024\_06 – Score 1.00 – global\_south\_2024\_06/109613190.txt}  \label{ex:f6_neg_269}
US intelligence has indicated that a large-scale confrontation between Israel and Hezbollah might break out in the upcoming weeks, if Jerusalem and Hamas fail to reach a ceasefire deal, media reports said.
US officials are trying to convince both sides to de-escalate -- a task that would be significantly easier with a cease-fire in place in Gaza, reported Politico.
The Israel Defense Forces and Hezbollah have drafted battle plans and are in the process of trying to procure additional weapons, according to two senior US officials briefed on the intelligence, reported Politico.
Both Israel and Hezbollah have repeatedly said they do not want to engage themselves in war.
The risk is higher now than at any other point in recent weeks, according to another senior US official, as quoted by Politico.
Meanwhile, US Secretary of Defense Lloyd J.
Austin III recently urged Israeli Defense Minister Yoav Gallant to work to de-escalate tensions with the Iranian-backed Hezbollah in Lebanon, during a meeting at the Pentagon on June 25, Pentagon Press Secretary Air Force Maj.
Gen.
Pat Ryder said. @ @ @ @ @ @ @ @ @ @ the Israel-Lebanon border, to surge more humanitarian aid into Gaza, and to stand together against Iranian and Iranian-supported attacks against Israel and destabilizing activities throughout the Middle East", Ryder said.
Hezbollah has reportedly escalated provocations on Northern Israel since the Hamas attack on Israel in October 7, 2023, and those provocations have been increasing.
Austin stressed that these provocations"threaten to drag the Israeli and Lebanese people into a war that neither of them wants and that such a war would be catastrophic for Lebanon and it would be devastating for innocent Israeli and Lebanese civilians", Ryder said as quoted by Pentagon website.
Austin told Gallant that diplomacy is the only way that can prevent further escalation of tensions in the region.
Meanwhile, delivery of humanitarian aid in Gaza remains a difficult and frustrating task for UN and partner agencies as the conflict there grinds on amid intense civilian suffering.
On Thursday, UN Spokesperson St? phane Dujarric, provided details from aid coordination office OCHA on how difficult @ @ @ @ @ @ @ @ @ @ to access constraints, food security and deteriorating health issues."Our partners report that limited access in the north is preventing the establishment of a new nutrition services centre in that area,"he told correspondents at UN Headquarters in New York.
OCHA is also struggling with finding enough space to establish nutrition sites in Deir al Balah and Khan Younis.
However, he said, the UN ' s partners are working to further their efforts in those areas as well as in Al Mawasi and Gaza City.

% matched lemmas: 
\end{textsample}
